\section{采样理论基础}

在讨论如何学习复杂的数据分布之前，授课教师首先回顾了一个基本问题：\textbf{如果我们已经知道了一个概率分布，如何从中采样？}

\subsection{逆变换采样（Inverse Transform Sampling）}

\subsubsection{离散情形：从 Categorical 分布采样}

考虑一个有 $K$ 个状态的 categorical 分布，每个状态 $k$ 有概率 $p_k$，且 $\sum_{k=1}^K p_k = 1$。从这个分布采样的经典方法是\textbf{逆变换采样}：

\begin{enumerate}
    \item 计算累积分布函数（CDF）：$F(k) = \sum_{i=1}^{k} p_i$
    \item 从均匀分布 $U \sim \text{Uniform}(0, 1)$ 中采样一个值 $u$
    \item 找到满足 $F(k-1) < u \leq F(k)$ 的状态 $k$，该状态即为采样结果
\end{enumerate}

直觉上，概率越高的状态在 CDF 的 $y$ 轴上占据的``区间''越宽，因此均匀采样落入该区间的概率也越大。

\subsubsection{连续情形}

对于连续随机变量 $X$，其概率密度函数为 $f(x)$，CDF 为：

$$F(x) = \int_{-\infty}^{x} f(t) \, dt$$

逆变换采样的步骤为：

\begin{enumerate}
    \item 采样 $u \sim \text{Uniform}(0, 1)$
    \item 计算 $x = F^{-1}(u)$
\end{enumerate}

\begin{importantbox}{逆变换采样的核心公式}
设 $F$ 为连续分布的 CDF，$F^{-1}$ 为其逆函数，则：
$$X = F^{-1}(U), \quad U \sim \text{Uniform}(0,1) \implies X \sim F$$
关键前提：\textbf{必须能够计算 CDF 的逆函数 $F^{-1}$}。
\end{importantbox}

\subsubsection{课堂练习：从 $f(x) = \frac{3}{8}x^2$ 采样}

授课教师给出了一个小练习：设 PDF 为 $f(x) = \frac{3}{8}x^2$，$x \in {[}-2, 2{]}$，求逆变换采样的采样公式。

解法：
\begin{enumerate}
    \item 计算 CDF：$F(x) = \int_{-2}^{x} \frac{3}{8}t^2 \, dt = \frac{1}{8}(x^3 + 8)$
    \item 令 $u = F(x)$，解出 $x$：$u = \frac{1}{8}(x^3 + 8) \implies x = (8u - 8)^{1/3} = 2(u-1)^{1/3}$
    \item 不对，更仔细的推导：$x^3 = 8u - 8 \implies x = \sqrt{[}3{]}{8u - 8}$
\end{enumerate}

因此采样过程为：先采 $u \sim \text{Uniform}(0,1)$，再计算 $x = (8u - 8)^{1/3}$。

\subsection{高斯分布的采样：Box-Muller 变换}

\subsubsection{问题的提出}

标准正态分布 $\mathcal{N}(0, 1)$ 的 PDF 为：

$$f(x) = \frac{1}{\sqrt{2\pi}} e^{-x^2/2}$$

其 CDF $\Phi(x) = \int_{-\infty}^{x} f(t) \, dt$ \textbf{没有解析形式的逆函数}，因此无法直接使用逆变换采样。

\begin{warningbox}{高斯 CDF 不可逆的陷阱}
虽然高斯分布的 PDF 有优美的解析形式，但其 CDF 的逆函数（即 probit 函数）不存在封闭形式。这意味着不能直接对高斯分布使用最简单的逆变换采样方法。实践中需要借助变量替换等技巧。
\end{warningbox}

\subsubsection{Box-Muller 变换的推导}

Box-Muller 变换的核心思想是：\textbf{同时考虑两个独立的标准正态样本，通过极坐标变换使得每个分量都可以用逆变换采样}。

设 $X, Y \overset{\text{i.i.d.}}{\sim} \mathcal{N}(0, 1)$，其联合密度为：

$$f_{X,Y}(x, y) = \frac{1}{2\pi} e^{-(x^2 + y^2)/2}$$

引入极坐标变量 $D = X^2 + Y^2$（平方半径）和 $\Theta$（角度），可以证明：

\begin{itemize}
    \item $\Theta \sim \text{Uniform}(0, 2\pi)$（角度均匀分布，直接采样）
    \item $D \sim \text{Exponential}(1/2)$，即 $f_D(d) = \frac{1}{2}e^{-d/2}$（$\chi^2$ 分布，自由度为 2）
\end{itemize}

对 $D$ 的分布，CDF 为 $F_D(d) = 1 - e^{-d/2}$，其逆函数为：

$$d = F_D^{-1}(u) = -2\ln(1 - u)$$

由于 $1-u$ 与 $u$ 同为 $\text{Uniform}(0,1)$，可简化为 $d = -2\ln(u)$。

\begin{importantbox}{Box-Muller 变换的完整公式}
给定两个独立的均匀分布样本 $U_1, U_2 \sim \text{Uniform}(0,1)$：
$$X = \sqrt{-2\ln U_1} \cdot \cos(2\pi U_2)$$
$$Y = \sqrt{-2\ln U_1} \cdot \sin(2\pi U_2)$$
则 $X, Y$ 为两个独立的标准正态分布样本。一次变换可以同时得到两个高斯样本。
\end{importantbox}

\subsection{多元高斯分布的采样与 Reparameterization}

\subsubsection{各向同性（Isotropic）高斯分布}

对于多元标准正态分布 $\mathcal{N}(\mathbf{0}, \mathbf{I})$，由于各维度独立，可以对每个维度分别用 Box-Muller 变换采样。

\subsubsection{各向异性（Anisotropic）高斯分布}

对于一般的多元高斯分布 $\mathcal{N}(\boldsymbol{\mu}, \boldsymbol{\Sigma})$，采样方法为\textbf{重参数化（Reparameterization）}：

$$\mathbf{x} = \boldsymbol{\mu} + \mathbf{L} \boldsymbol{\epsilon}, \quad \boldsymbol{\epsilon} \sim \mathcal{N}(\mathbf{0}, \mathbf{I})$$

其中 $\mathbf{L}$ 满足 $\boldsymbol{\Sigma} = \mathbf{L}\mathbf{L}^T$（例如 Cholesky 分解）。

\begin{importantbox}{Reparameterization Trick}
从 $\mathcal{N}(\boldsymbol{\mu}, \boldsymbol{\Sigma})$ 采样等价于：
$$\mathbf{x} = \boldsymbol{\mu} + \boldsymbol{\Sigma}^{1/2} \boldsymbol{\epsilon}, \quad \boldsymbol{\epsilon} \sim \mathcal{N}(\mathbf{0}, \mathbf{I})$$
这个技巧不仅是采样的基础，更是 VAE 中实现梯度反向传播的关键——它将随机性``分离''到 $\boldsymbol{\epsilon}$ 中，使得采样过程对 $\boldsymbol{\mu}$ 和 $\boldsymbol{\Sigma}$ 可微。
\end{importantbox}

\begin{knowledgebox}{为什么生成模型偏爱高斯分布？}
高斯分布在生成模型中无处不在，原因包括：（1）我们知道如何高效采样；（2）高斯分布具有优良的数学性质（如 KL 散度有闭合形式）；（3）中心极限定理保证许多分布渐近趋向高斯；（4）在给定均值和方差的约束下，高斯分布是最大熵分布。
\end{knowledgebox}

\subsection{本章小结}

当概率分布的形式已知时，我们可以通过逆变换采样、变量替换（如 Box-Muller 变换）以及重参数化等技术进行采样。然而，现实中的数据分布（如自然图像的分布）形式未知，且极其复杂——这正是需要深度生成模型的原因。

\newpage
